\documentclass[a4paper,12pt]{article}

% Pacotes básicos
\usepackage[utf8]{inputenc}
\usepackage[T1]{fontenc}
\usepackage[brazil]{babel}
\usepackage{geometry}
\usepackage{setspace}
\usepackage{hyperref}
\usepackage{listings}
\usepackage{xcolor}
\usepackage{longtable}
\usepackage{changepage}

% Configuração da página
\geometry{margin=2.5cm}
\onehalfspacing

% Configuração de código
\lstset{
    basicstyle=\ttfamily\small,
    keywordstyle=\color{blue},
    commentstyle=\color{gray},
    stringstyle=\color{orange},
    breaklines=true,
    frame=single,
    showspaces=false,
    showstringspaces=false
}

% Syntax highlighting da linguagem personalizada
\lstdefinelanguage{Dhara}{
    % keywords by group — each group can have its own color
    morekeywords=[1]{style, borderline, PREY, HOMETOWN},
    morekeywords=[2]{space, lithium, judas},
    morekeywords=[3]{houdini, more, problems, bloomfield, loser},
    morekeywords=[4]{catapult, pleaser},
    sensitive=true,        % case-sensitive keywords
    morestring=[b]",       % strings delimited by "
    morecomment=[l]{~~},   % line comments starting with ~~
}

\lstdefinestyle{DharaStyle}{
    language=Dhara,
    basicstyle=\ttfamily\small,
    keywordstyle=[1]\color{violet}\bfseries,   % style, borderline, PREY, HOMETOWN
    keywordstyle=[2]\color{blue}\bfseries,     % space, lithium, judas
    keywordstyle=[3]\color{orange}\bfseries,   % houdini, more, problems, bloomfield
    keywordstyle=[4]\color{teal}\bfseries,     % catapult, pleaser
    stringstyle=\color{red},
    commentstyle=\color{gray}\itshape,
    showstringspaces=false,
    frame=single,
    breaklines=true,
}

\lstdefinelanguage{Rust}{
    keywords={
        fn, let, mut, pub, impl, trait, struct, enum, match, if, else,
        while, for, loop, return, break, continue, use, mod, crate,
        self, Self, super, where, as, const, static, ref, type
    },
    keywordstyle=\color{blue}\bfseries,
    ndkeywords={i32, i64, u32, u64, f32, f64, bool, String, Vec},
    ndkeywordstyle=\color{teal},
    identifierstyle=\color{black},
    sensitive=true,
    comment=[l]{//},
    morecomment=[s]{/*}{*/},
    commentstyle=\color{gray}\ttfamily,
    stringstyle=\color{orange},
    morestring=[b]",
}

\title{\textbf{Linguagem DPP}}
\author{
    Gabriel Santana Dias --- R.A. 22.225.017-7\\
    Pedro Schneider \phantom{Dias} --- R.A. 22.225.016-9
}
\date{\today}

\begin{document}

\maketitle

\tableofcontents
\newpage

% ==========================
\section{Introdução}
Este projeto implementa uma linguagem fictícia chamada \textbf{Dhara} (DPP), cujas palavras reservadas são baseadas em títulos de músicas.

% ==========================
\section{Analisador Léxico}
\subsection{Expressões Regulares}

\begin{itemize}

    \item \textbf{Palavra}:  
    \texttt{[a-zA-Z]+[a-zA-Z]*}

    \item \textbf{Número}:  
    \texttt{[0-9]+}

    \begin{itemize}
        \item \textbf{Identificador}:  
        \texttt{palavra+num*}
        
        \item \textbf{Número inteiro}:  
        \texttt{num+}

        \item \textbf{Número float}:  
        \texttt{num+ '.' num+}
        
        \item \textbf{String}:  
        \texttt{" palavra* "}
        
        \item \textbf{Comentário de linha}:
        \texttt{`\~{}\~{}` palavra* `\textbackslash n`}
    \end{itemize}
    
    \item \textbf{Palavras reservadas}:
        \begin{lstlisting}[basicstyle=\ttfamily\footnotesize]
         style   | borderline |
         space   |  lithium   | judas
        pleaser  |  catapult  |
        houdini  |    more    |
        problems | bloomfield |  not...ok
          PREY   |  HOMETOWN  |
        loser
        \end{lstlisting}

    \item \textbf{Operadores lógicos}:
    \texttt{\&{}\&{} \ \ ||}
        
    \item \textbf{Operadores comparativos}:
    \texttt{== \ \ !=  \ \ >= \ \ <= \ \ > \ \ <}

    \item \textbf{Operadores aritméticos}:
    \texttt{+ \ \ - \ \ * \ \ / \ \ \%}

    \item \textbf{Delimitadores}:
    \texttt{( \ ) \ \{ \ \}}

    \item \textbf{Quebra de linha}:
    \texttt{\textbackslash n}

\end{itemize}

% ==========================
\newpage
\section{Analisador Semântico}
\subsection{Gramática Livre de Contexto (GLC)}

\begin{adjustwidth}{-2cm}{-2cm}
\begin{lstlisting}[basicstyle=\ttfamily\scriptsize]
palavra    -> [a-zA-Z]+[a-zA-Z]*
num        -> [0-9]+

id         -> palavra+num*
int        -> num+
float      -> num+ '.' num+
string     -> " palavra* "

op_logic   -> && | ||
op_comp    -> == | != | >= | <= | > | <
op_arit    ->  + |  - |  * |  / | %

tipo       -> `space` | `lithium` | `judas`
atribuicao -> tipo id = expressao ; | id = expressao ;
declaracao -> tipo id ;

comentario -> `~~` palavra* `\n`

expressao  -> termo expressao'
expressao' -> op_arit termo expressao' | e
termo      -> ( expressao ) | int | float | string | id | chamada_funcao

condicao   -> termo condicao'
condicao'  -> (op_logic | op_comp) termo condicao' | e

saida             -> `catapult` [ expressao ] ;

entrada           -> `pleaser` [ entrada' ] ;
entrada'          -> " ponteiros " , identificadores
ponteiros         -> formatadores ponteiros'
ponteiros'        -> , ponteiros | e
identificadores   -> id identificadores'
identificadores'  -> , identificadores | e
formatadores      -> %d | %f | %s

declaracao_funcao       -> tipo `PREY` id [ parametros_declaracao ] { comando* `HOMETOWN` expressao? }
parametros_declaracao   -> tipo id parametros_declaracao'
parametros_declaracao'  -> , tipo id parametros_declaracao' | e

chamada_funcao          -> id [ parametros_chamada ] ;
parametros_chamada      -> expressao parametros_chamada'
parametros_chamada'     -> , expressao parametros_chamada' | e

if       -> `houdini` ( condicao ) { comando* } else
else     -> `more` { comando* } | e
while    -> `problems` ( condicao ) { comando* }
do_while -> `not...ok` { comando* } while ( condicao ) ;
for      -> `bloomfield` ( atribuicao condicao atribuicao ) { comando* }

comando  -> comentario | atribuicao | entrada | saida | if | while | do_while | for | chamada_funcao
main     -> declaracao_funcao* `style` codigo `borderline`
codigo   -> comando* codigo* | e
\end{lstlisting}
\end{adjustwidth}

% ==========================
% ==========================
\section{Execução do Compilador}

O compilador foi desenvolvido em Java e está organizado na seguinte estrutura de diretórios:

\begin{lstlisting}
src/
    analisadorLexico/
    analisadorSintatico/
    Main.java
    ParserTest.java
\end{lstlisting}

Onde \texttt{ParserTest.java} contém testes unitários do analisador sintático. 

\subsection{Pré-requisitos}
\begin{itemize}
    \item Java JDK 17 ou superior
    \item Código na linguagem dpp num arquivo .txt em \texttt{src/}
\end{itemize}

\subsection{Compilação}

A compilação deve ser feita a partir da raiz do projeto:

\begin{lstlisting}
javac -d bin src/*.java src/analisadorLexico/*.java src/analisadorSintatico/*.java
\end{lstlisting}

O comando acima gera os arquivos compilados no diretório \texttt{bin/}.

\subsection{Execução}

Após a compilação, execute o compilador com:

\begin{lstlisting}
java -cp bin Main codigo.txt -o output.rs
\end{lstlisting}

Use as flags `-t` para ver a lista de tokens (\emph{lex table}) ou `-a` para ver a AST (árvore de derivação da GLC).  
Use a flag `-o` para definir o nome do arquivo .rs de saída (código escrito  em rust)

% ==========================

\section{Características da Linguagem}

A linguagem \textbf{Dhara (DPP)} foi projetada com sintaxe inspirada em linguagens tradicionais, como C.

\begin{itemize}
    \item \textbf{Tipagem}: linguagem de tipagem estática e forte, exigindo declaração explícita dos tipos das variáveis.
    
    \item \textbf{Tipos de dados suportados}:
    \begin{itemize}
        \item \texttt{space} --- números inteiros
        \item \texttt{lithium} --- números de ponto flutuante
        \item \texttt{judas} --- strings
    \end{itemize}

    \item \textbf{Estruturas condicionais}:
    \begin{itemize}
        \item \texttt{houdini} --- equivalente ao \texttt{if}
        \item \texttt{more} --- equivalente ao \texttt{else}
    \end{itemize}

    \item \textbf{Estruturas de repetição}:
    \begin{itemize}
        \item \texttt{problems} --- equivalente ao \texttt{while}
        \item \texttt{bloomfield} --- equivalente ao \texttt{for}
        \item \texttt{not...ok} --- equivalente ao \texttt{do while}
    \end{itemize}

    \item \textbf{Funções}:
    \begin{itemize}
        \item Declaração de funções através da palavra reservada \texttt{PREY}
        \item Retorno de valores utilizando \texttt{HOMETOWN}
        \item Suporte a parâmetros e chamadas de função
    \end{itemize}

    \item \textbf{Operadores suportados}:
    \begin{itemize}
        \item Aritméticos: \texttt{+, -, *, /, \%}
        \item Comparação: \texttt{==, !=, >, <, >=, <=}
        \item Lógicos: \texttt{\&\&, ||}
    \end{itemize}

    \item \textbf{Entrada e saída}:
    \begin{itemize}
        \item \texttt{catapult} --- saída de dados no terminal
        \item \texttt{pleaser} --- entrada de dados pelo usuário
    \end{itemize}

    \item \textbf{Comentários}:
    \begin{itemize}
        \item Comentários de linha iniciados por \texttt{\~{}\~{}}
    \end{itemize}

    \item \textbf{Estrutura principal do programa}:
    \begin{itemize}
        \item O programa principal inicia com \texttt{style} e termina com \texttt{borderline}
    \end{itemize}
\end{itemize}

% ==========================
\newpage
\section{Exemplos de Código}

\subsection{Código na Linguagem Criada}
\begin{adjustwidth}{-1cm}{-1cm}
\begin{lstlisting}[
    style=DharaStyle,
    basicstyle=\ttfamily\footnotesize,
    numbers=left,
    numberstyle=\tiny\color{gray},
    stepnumber=1,
    numbersep=8pt
]
lithium PREY multiplicar[lithium a, lithium b] {
    ~~ multiplica dois numeros
    HOMETOWN a*b;
}

style
space x;
space q;
lithium y;

bloomfield(x=0; x<=10; x++) {
    catapult[x];
}

houdini((x == 10) && (x != 0)) {
    catapult["Digite um numero inteiro e um decimal: "];
    pleaser["%d,%f", q, y];
    catapult["Voce digitou: " + q + " e " + y];
}

houdini(y < 10.5) {
    problems(y < 11) {
        y++;
    }
} more {
    lithium res = multiplicar[q, y];
    catapult["Q*Y e igual a: " + res];
}
borderline
\end{lstlisting}
\end{adjustwidth}

\newpage
\subsection{Código Traduzido (Rust)}
\begin{adjustwidth}{-1cm}{-1cm}
\begin{lstlisting}[
    language=Rust,
    basicstyle=\ttfamily\footnotesize,
    numbers=left,
    numberstyle=\tiny\color{gray},
    stepnumber=1,
    numbersep=8pt,
    breaklines=true,
    frame=single
]
fn multiplicar(a: f32, b: f32) -> f32 {
    // multiplica dois numeros
    a * b
}

fn main() {
    let mut x: i32;
    let mut q: i32 = 0;
    let mut y: f32 = 0.0;

    for i in 0..=10 {
        x = i;
        println!("{}", x);
    }

    x = 10;

    if (x == 10) && (x != 0) {
        use std::io;

        println!("Digite um numero inteiro e um decimal:");

        let mut entrada = String::new();
        io::stdin().read_line(&mut entrada).unwrap();

        let valores: Vec<&str> =
            entrada.trim().split_whitespace().collect();

        q = valores[0].parse().unwrap();
        y = valores[1].parse().unwrap();

        println!("Voce digitou: {} e {}", q, y);
    }

    if y < 10.5 {
        while y < 11.0 {
            y += 1.0;
        }
    } else {
        let res: f32 = multiplicar(q as f32, y);
        println!("Q*Y e igual a: {}", res);
    }
}
\end{lstlisting}
\end{adjustwidth}

% ==========================

\end{document}
